\chapter{Options-Implied Signals for Stocks}\label{ch:s1:options-implied-signals-for-stocks}

When a stock's out-of-the-money puts become expensive relative to its at-the-money options, someone is paying up for protection, and the stock tends to fall.
Xing, Zhang and Zhao found that stocks with the steepest volatility smirks underperformed those with the flattest by 10.9\% a year, risk-adjusted, and had
the worst earnings surprises the following quarter: informed traders with bad news prefer puts. The option market is where leverage is cheap and shorting is
unnecessary, so it is where some information shows first. This chapter adds an options layer to the synthetic market, with informed traders who buy options
ahead of planted earnings news, and measures what skew, the call--put volatility spread and option volume predict, how fast it decays, and why it is best
traded around events. The build is \texttt{firm.optsignal}.

\section{Skew as a signal}

\begin{definition}[Informed options trading]\label{def:s1:options-implied-signals-for-stocks:informed}
\emph{Informed options trading}\index{informed options trading} is trading in a stock's options by investors with private information about the stock, who
choose options over the stock for their leverage and, for bad news, to avoid the costs and constraints of short selling; their demand moves the prices and
volumes of the options they buy before the information reaches the stock.
\end{definition}

The volatility skew (Book~1, chapter~25) is the extra implied volatility of out-of-the-money puts over at-the-money options. It is there on every stock (a
premium for crash protection), and it moves when buyers of puts outnumber sellers. In the synthetic options layer every stock has an at-the-money implied
volatility equal to its realised volatility plus a two-point premium, a put skew of three points, and daily noise of one point in each; informed traders,
in the ten days before each earnings announcement (chapter~7's calendar), buy puts ahead of bad surprises and calls ahead of good ones, and each
standard deviation of surprise moves the implied volatility of the options they buy by 0.3 points. Bad news is weighted 1.5 times good news, for the reason
Johnson and So give below. The skew's sign as a signal is negative: a steep skew predicts a fall.

\section{Implied minus realised volatility}

\begin{definition}[Implied volatility spread]\label{def:s1:options-implied-signals-for-stocks:spread}
A stock's \emph{implied volatility spread}\index{implied volatility spread} is the difference between the implied volatilities of its calls and its puts at
matched strikes and maturities; under put--call parity it is zero, and a positive spread means calls are relatively expensive.
\end{definition}

Cremers and Weinbaum measured deviations from put--call parity this way and found that stocks with relatively expensive calls outperformed those with
relatively expensive puts by 50 basis points a week, not explained by short-sale constraints, stronger where options were liquid and stocks were not, and
declining over their sample. Bali and Hovakimian separated two spreads: the call--put spread (positively related to future returns, a proxy for jump
risk and informed trading) and the implied--realised spread (negatively related to returns through realised minus implied, a proxy for volatility risk);
levels of volatility predicted nothing. In the synthetic layer the implied--realised spread carries no planted information, and its IC is zero at every
horizon: a useful control.

\begin{center}
\small
\begin{tabular}{@{}lcccc@{}}
\toprule
mean daily rank IC, years 3 to 10 & 1 day & 5 days & 20 days & 60 days\\
\midrule
skew (minus), all stocks & 0.004 & 0.014 & 0.018 & 0.012\\
call--put volatility spread, all stocks & 0.005 & 0.017 & 0.022 & 0.015\\
option-to-stock volume (minus), all stocks & 0.002 & 0.008 & 0.010 & 0.007\\
implied minus realised volatility, all stocks & 0.000 & 0.001 & 0.000 & $-0.000$\\
\midrule
skew (minus), announcing within ten days & 0.026 & 0.085 & 0.111 & 0.082\\
call--put volatility spread, announcing within ten days & 0.030 & 0.101 & 0.129 & 0.094\\
option-to-stock volume (minus), announcing within ten days & 0.015 & 0.047 & 0.061 & 0.045\\
\bottomrule
\end{tabular}
\end{center}

\section{Option volume and the put--call ratio}

\begin{definition}[Put--call ratio, option-to-stock volume ratio]\label{def:s1:options-implied-signals-for-stocks:volume}
A stock's \emph{put--call ratio}\index{put--call ratio} is its put option volume divided by its call option volume over a period, ideally counting only
volume that opens new long positions. Its \emph{option-to-stock volume ratio}\index{option-to-stock volume ratio} is its total option volume (in shares of the
underlying) divided by its stock volume.
\end{definition}

Pan and Poteshman built \omterm{def:s1:options-implied-signals-for-stocks:volume}{put--call ratios} from buyer-initiated volume that opened new positions, a data set most researchers do not have, and found that
stocks with low ratios beat those with high ratios by more than 40 basis points the next day and more than 1\% over the next week, from nonpublic
information held by option traders. Johnson and So showed why unsigned volume is informative too: short-sale costs make informed traders use options more
for bad news than for good, so a high \omterm{def:s1:options-implied-signals-for-stocks:volume}{option-to-stock volume ratio} signals bad news; the lowest decile beat the highest by 0.34\% a week. In the synthetic
layer option volume rises with the size of the informed demand whatever its direction, weighted toward bad news, so its IC is about half that of the price
signals: volume says someone knows something, and the asymmetry says what.

\section{Where informed traders trade}

The table's lower half is the chapter's lesson. Across all stocks the signals' ICs are small (0.02 for the volatility spread at twenty days), because most
stocks on most days have no informed traders in their options and the signals are noise. Among stocks announcing within ten days (11.7\% of stock-days) the
same signals have ICs about six times larger: the information is concentrated where the events are. The ICs rise to twenty days and then decay, since the
announcement's jump, which the informed traders anticipated, falls inside the twenty-day window and the planted post-announcement drift carries part of it
further.

\begin{omfigure}
\begin{tikzpicture}
\begin{axis}[width=10.5cm,height=5.4cm,xmode=log,xlabel={\small horizon (days)},ylabel={\small mean daily rank IC},tick label style={font=\footnotesize},
  xtick={1,5,20,60},xticklabels={1,5,20,60},xmin=0.8,xmax=75,ymin=0,ymax=0.14,ytick={0,0.05,0.1},
  yticklabel style={/pgf/number format/fixed,/pgf/number format/precision=2},grid=major,grid style={gray!25},
  legend style={font=\scriptsize,at={(1.02,1)},anchor=north west},table/col sep=comma]
\addplot[omDef,very thick,mark=*] table[x=h,y=spread_ev]{figdata/strategies-1/15-options-implied-signals-for-stocks/ic.csv};
\addplot[omThm,very thick,mark=square*] table[x=h,y=skew_ev]{figdata/strategies-1/15-options-implied-signals-for-stocks/ic.csv};
\addplot[black!50,thick,mark=triangle*] table[x=h,y=os_ev]{figdata/strategies-1/15-options-implied-signals-for-stocks/ic.csv};
\addplot[omDef,thick,dashed,mark=o] table[x=h,y=spread]{figdata/strategies-1/15-options-implied-signals-for-stocks/ic.csv};
\addplot[omThm,thick,dashed,mark=square] table[x=h,y=skew]{figdata/strategies-1/15-options-implied-signals-for-stocks/ic.csv};
\legend{spread (events),skew (events),volume (events),spread (all),skew (all)}
\end{axis}
\end{tikzpicture}

\omcaption{Rank IC of the options signals against the market-adjusted return over the next 1 to 60 days on the synthetic market with informed options
traders: among stocks announcing within ten days (solid) and all stocks (dashed). Data: \texttt{s1\_options.ic}.}
\label{fig:s1:options-implied-signals-for-stocks:ic}
\end{omfigure}

The books follow. A weekly decile book on the volatility spread across all stocks earns 8.5\% a year before costs (Sharpe ratio 3.67) and loses 0.9\% after ten
basis points per unit traded: the implied volatilities' daily noise churns the ranking. Restricted to stocks announcing within ten days (quintiles of that
smaller set), it earns 31.9\% a year after costs at a Sharpe ratio of 6.12, the skew book 23.0\% and the volume book 9.1\%. Those synthetic numbers are far
above anything in the literature, and the reason is the model: its informed traders know the whole surprise and their demand is a clean function of it. Real
informed flow is a small, unsigned part of option volume, shared with hedgers and speculators; what carries over is the structure: trade the signals where
the events are, and rebalance only as fast as the information changes.

\section{Strategy files}

\begin{strategyfile}{Skew signal}\label{strat:s1:options-implied-signals-for-stocks:skew}
\sfield{Who pays you, and why} Informed traders who reveal bad news by buying puts, and the stock market's slowness to follow.
\sfield{Instruments and venues} Stocks with liquid options; the signal from the options, the trade in the stock.
\sfield{Signal} Out-of-the-money put implied volatility minus at-the-money, against the stock's own history.
\sfield{Sizing and execution} Short steep skews, long flat ones; concentrate around scheduled events.
\sfield{Costs} Stock trading; the signal's noise drives turnover.
\sfield{How it dies} Crowding; skew driven by hedging demand rather than information.
\sfield{Horizon, capacity, infrastructure} Weeks to months; option surfaces by stock (Book~5).
\sfield{Backtest honestly} Implied volatilities from closing quotes before the stock trade, with their bid--ask; the event calendar known in advance.
\sfield{Sources} Xing, Zhang and Zhao (2010): 10.9\% a year risk-adjusted between extreme smirk portfolios.
\end{strategyfile}

\begin{strategyfile}{Implied-volatility spread signal}\label{strat:s1:options-implied-signals-for-stocks:spread}
\sfield{Who pays you, and why} As for skew, from both sides: expensive calls signal good news, expensive puts bad.
\sfield{Instruments and venues} Stocks with liquid options at matched strikes.
\sfield{Signal} Call minus put implied volatility, open-interest weighted across strikes.
\sfield{Sizing and execution} Long high spreads, short low; weekly or around events.
\sfield{Costs} As for skew.
\sfield{How it dies} It weakened over Cremers and Weinbaum's sample.
\sfield{Horizon, capacity, infrastructure} Days to weeks.
\sfield{Backtest honestly} Early-exercise premia removed from American options before computing the spread.
\sfield{Sources} Cremers and Weinbaum (2010): 50 basis points a week; Bali and Hovakimian (2009).
\end{strategyfile}

\begin{strategyfile}{Option volume signal}\label{strat:s1:options-implied-signals-for-stocks:volume}
\sfield{Who pays you, and why} Informed traders whose option volume rises before news, more so for bad news.
\sfield{Instruments and venues} Stocks with option volume data.
\sfield{Signal} The \omterm{def:s1:options-implied-signals-for-stocks:volume}{option-to-stock volume ratio} (unsigned) or a \omterm{def:s1:options-implied-signals-for-stocks:volume}{put--call ratio} of opening buys (signed).
\sfield{Sizing and execution} Short high ratios; with signed data, long low \omterm{def:s1:options-implied-signals-for-stocks:volume}{put--call ratios}.
\sfield{Costs} As above.
\sfield{How it dies} Volume from hedging and structured products that swamps informed flow.
\sfield{Horizon, capacity, infrastructure} Days to a week; volume by option series.
\sfield{Backtest honestly} Volume as reported that day, not revised; signed volume only where the data really is signed.
\sfield{Sources} Pan and Poteshman (2006): more than 1\% a week; Johnson and So (2012): 0.34\% a week between extreme deciles.
\end{strategyfile}

\begin{strategyfile}{Implied-minus-realised volatility signal}\label{strat:s1:options-implied-signals-for-stocks:ivrv}
\sfield{Who pays you, and why} Investors paying for volatility insurance, and the risk premium that insurance carries.
\sfield{Instruments and venues} Stocks and their options.
\sfield{Signal} Realised minus implied volatility (Bali and Hovakimian's sign).
\sfield{Sizing and execution} Cross-sectional sorts, monthly.
\sfield{Costs} Low turnover.
\sfield{How it dies} As a risk premium, in volatility spikes.
\sfield{Horizon, capacity, infrastructure} Months; volatility surfaces.
\sfield{Backtest honestly} Realised volatility from data before the implied quote.
\sfield{Sources} Bali and Hovakimian (2009); this chapter's control (IC zero where nothing is planted).
\end{strategyfile}

\section{Tutorial: someone is buying protection}\label{tut:s1:options-implied-signals-for-stocks:tut}

\begin{tutorial}
\textbf{Goal.} Add an options layer with informed traders to the synthetic market, compute the options signals, and measure them by horizon and around
events. \textbf{End state:} the table and \cref{fig:s1:options-implied-signals-for-stocks:ic}.
\begin{tutsteps}
\item \textbf{The layer}: implied volatilities and volume, moved by informed demand before announcements.
\omcode{code/firm/optsignal/firm_optsignal.py}{37}{58}{Informed options demand and the signals it leaves.}\label{lst:s1:options-implied-signals-for-stocks:layer}
\item \textbf{Run} \texttt{ic(name, h)} and \texttt{ic(name, h, True)} for the four signals and four horizons, \texttt{book(name)} and \texttt{book(name,
True)}, and \texttt{fig\_options.py}.
\end{tutsteps}
\textbf{What to change next.} Let most option volume be hedging and noise, with informed flow a small share; smooth the signals to cut turnover; add a
cost of shorting and let informed traders choose between the stock and the options.
\end{tutorial}

\section{Build: the options layer}\label{bld:s1:options-implied-signals-for-stocks:optsignal}

\begin{build}
\bfield{Purpose} A synthetic option market attached to the stock panel, and the option-implied signals built from it.
\bfield{Interface} \texttt{OptionConfig}, \texttt{simulate\_options(flag, surprise, listed, rv, cfg, rng)}, \texttt{signals(opt, rv)}.
\bfield{Rules} Informed demand only before events and in the direction of their surprise; everything else noise; signals from the day's closing values.
\bfield{Acceptance tests} \texttt{code/firm/optsignal/tests/}: the right options move, by the right amount, only in the window before the event; volume rises
more for bad news; implied minus realised equals the premium without noise.
\bfield{Stretch} A full volatility surface per stock (Book~5); hedging and structured-product flow; signed volume.
\end{build}

\begin{omsources}
\item Y.~Xing, X.~Zhang and R.~Zhao, ``What does the individual option volatility smirk tell us about future equity returns?'', \emph{Journal of Financial and
Quantitative Analysis} 45(3), 2010.
\item M.~Cremers and D.~Weinbaum, ``Deviations from put-call parity and stock return predictability'', \emph{Journal of Financial and Quantitative Analysis}
45(2), 2010.
\item J.~Pan and A.~M.~Poteshman, ``The information in option volume for future stock prices'', \emph{Review of Financial Studies} 19(3), 2006.
\item T.~L.~Johnson and E.~C.~So, ``The option to stock volume ratio and future returns'', \emph{Journal of Financial Economics} 106(2), 2012.
\item T.~G.~Bali and A.~Hovakimian, ``Volatility spreads and expected stock returns'', \emph{Management Science} 55(11), 2009.
\end{omsources}

\section{Exercises}

\begin{exercise}[$\star$]\label{exo:s1:options-implied-signals-for-stocks:1}
In the synthetic layer, how much does a surprise of $-2$ standard deviations raise the put implied volatility in the days before the announcement? And a
surprise of $+2$ the call implied volatility?
\end{exercise}

\begin{exercise}[$\star$]\label{exo:s1:options-implied-signals-for-stocks:2}
Johnson and So's spread of 0.34\% a week is how much a year, without compounding? And Cremers and Weinbaum's 50 basis points a week?
\end{exercise}

\begin{exercise}[$\star$]\label{exo:s1:options-implied-signals-for-stocks:3}
Why is the skew's sign as a signal negative and the call--put spread's positive?
\end{exercise}

\begin{exercise}[$\star\star$]\label{exo:s1:options-implied-signals-for-stocks:4}
Why might informed traders prefer options for bad news more than for good news?
\end{exercise}

\begin{exercise}[$\star\star$]\label{exo:s1:options-implied-signals-for-stocks:5}
The ICs are about six times larger among stocks announcing within ten days. Why, and what does it suggest for the book?
\end{exercise}

\begin{exercise}[$\star\star$]\label{exo:s1:options-implied-signals-for-stocks:6}
Why does the all-stock weekly book lose after costs when its gross Sharpe ratio is above 3?
\end{exercise}

\begin{exercise}[$\star\star\star$]\label{exo:s1:options-implied-signals-for-stocks:7}
\emph{Coding.} Run \texttt{book('os', True)} and compare it with the price-based books. Why is volume the weakest signal here?
\end{exercise}

\begin{exercise}[$\star\star\star$]\label{exo:s1:options-implied-signals-for-stocks:8}
\emph{Find the flaw.} ``The volatility-spread book earns 32\% a year after costs on the synthetic market; we should expect a large part of that in practice.''
\end{exercise}

\section{Problem: Someone Is Buying Protection}

\begin{problem}[{Weekend problem --- information in the options market}]\label{pb:s1:options-implied-signals-for-stocks:1}
The chapter's options layer and the public record.

\medskip\noindent\textbf{Part I --- The signals.}
\begin{enumerate}
\item Define \omterm{def:s1:options-implied-signals-for-stocks:informed}{informed options trading}, the \omterm{def:s1:options-implied-signals-for-stocks:spread}{implied volatility spread}, the \omterm{def:s1:options-implied-signals-for-stocks:volume}{put--call ratio} and the \omterm{def:s1:options-implied-signals-for-stocks:volume}{option-to-stock volume ratio}.
\item What did Xing, Zhang and Zhao find about the smirk?
\item What did Cremers and Weinbaum, and Bali and Hovakimian, find about spreads?
\item What did Pan and Poteshman, and Johnson and So, find about volume?
\end{enumerate}

\medskip\noindent\textbf{Part II --- The layer.}
\begin{enumerate}[resume]
\item Describe the synthetic options layer and its informed traders.
\item Why is bad news weighted more?
\item Why is implied minus realised a control here?
\item How is the event calendar used?
\end{enumerate}

\medskip\noindent\textbf{Part III --- The measurements.}
\begin{enumerate}[resume]
\item Give the ICs by horizon for all stocks.
\item Give them for stocks announcing within ten days.
\item Why do the ICs peak at twenty days?
\item Give the books' results, all stocks and around events.
\end{enumerate}

\medskip\noindent\textbf{Part IV --- The verdict.}
\begin{enumerate}[resume]
\item State the \emph{named result}: the IC of the skew and volatility-spread signals and their decay with horizon.
\item Why are the synthetic returns far above the literature's?
\item What would you change to make the layer realistic?
\item How would you cut the all-stock book's turnover?
\item What data would you need to trade these strategies?
\item Which signal decayed in the published record?
\item How do these signals relate to chapter~7's earnings strategies?
\item In one sentence: why look at options to trade stocks?
\end{enumerate}
\end{problem}

\section{Interview questions}

\begin{interviewq}[$\star$ \iqroles{researcher, trader}]\label{iq:s1:options-implied-signals-for-stocks:1}
Why might option prices predict stock returns?
\end{interviewq}

\begin{interviewq}[$\star\star$ \iqroles{researcher}]\label{iq:s1:options-implied-signals-for-stocks:2}
How would you compute a stock's volatility skew for a cross-sectional signal?
\end{interviewq}

\begin{interviewq}[$\star\star$ \iqroles{researcher}]\label{iq:s1:options-implied-signals-for-stocks:3}
What is the difference between the call--put \omterm{def:s1:options-implied-signals-for-stocks:spread}{implied volatility spread} and implied minus realised volatility as signals?
\end{interviewq}

\begin{interviewq}[$\star\star$ \iqroles{trader}]\label{iq:s1:options-implied-signals-for-stocks:4}
A stock's put volume triples ahead of earnings. What do you do?
\end{interviewq}

\begin{interviewq}[$\star\star$ \iqroles{risk}]\label{iq:s1:options-implied-signals-for-stocks:5}
What are the pitfalls of backtesting option-implied signals?
\end{interviewq}

\begin{interviewq}[$\star\star\star$ \iqroles{researcher}]\label{iq:s1:options-implied-signals-for-stocks:6}
In a model where informed traders choose between shorting the stock at a cost $c$ and buying puts with leverage $L$, derive when they prefer the puts, and
what that implies for the information in put volume.
\end{interviewq}
